%=======================================================================
%  Informe del Taller Práctico: Terraform + Docker
%
%  Documento de un solo fichero. Compilar con:
%     latexmk -pdf main.tex      (o simplemente:  make)
%=======================================================================
\documentclass[11pt,a4paper]{article}

%----------------------- Codificación e idioma -------------------------
\usepackage[T1]{fontenc}
\usepackage[utf8]{inputenc}
\usepackage[spanish,es-tabla,es-noquoting]{babel}
\usepackage{lmodern}
\usepackage{microtype}

%------------------------------ Geometría ------------------------------
\usepackage[a4paper,top=2.5cm,bottom=2.5cm,left=2.6cm,right=2.4cm,
            headheight=14pt]{geometry}

%--------------------------- Tablas y listas ---------------------------
\usepackage{booktabs}
\usepackage{tabularx}
\usepackage{ragged2e}
\newcolumntype{Y}{>{\RaggedRight\arraybackslash}X}
\renewcommand{\arraystretch}{1.3}
\usepackage{enumitem}
\setlist{itemsep=3pt,parsep=0pt,topsep=5pt}

%--------------------------- Figuras y color ---------------------------
\usepackage{graphicx}
\graphicspath{{figuras/}}
\usepackage{float}
\usepackage[table]{xcolor}
\usepackage{caption}
\captionsetup{font=small,labelfont=bf,skip=6pt}

%------------------------- Bloques de comandos -------------------------
\usepackage{listings}
\definecolor{acento}{HTML}{1168BD}
\definecolor{fondocodigo}{HTML}{F4F4F6}
\definecolor{bordecodigo}{HTML}{D7D7DE}
\lstdefinestyle{shellstyle}{
  basicstyle       = \ttfamily\small,
  backgroundcolor  = \color{fondocodigo},
  frame            = single,
  rulecolor        = \color{bordecodigo},
  framesep         = 6pt,
  xleftmargin      = 2pt,
  xrightmargin     = 2pt,
  breaklines       = true,
  columns          = fullflexible,
  keepspaces       = true,
  showstringspaces = false,
  aboveskip        = 10pt,
  belowskip        = 10pt}
\lstnewenvironment{shell}{\lstset{style=shellstyle}}{}

%--------------------------- Estilo de página --------------------------
\usepackage{fancyhdr}
\pagestyle{fancy}
\fancyhf{}
\fancyhead[L]{\small Taller Terraform + Docker}
\fancyhead[R]{\small\thepage}
\renewcommand{\headrulewidth}{0.4pt}

\usepackage{titlesec}
\titleformat{\section}{\normalfont\Large\bfseries\color{acento}}{\thesection.}{8pt}{}
\titleformat{\subsection}{\normalfont\large\bfseries}{\thesubsection}{8pt}{}

\usepackage[hidelinks]{hyperref}
\hypersetup{
  pdftitle  = {Informe del Taller Práctico: Terraform + Docker},
  pdfauthor = {Jairo Miguel Lara Serrano}}

% Atajo para una figura de pantalla a ancho completo.
\newcommand{\captura}[2]{%
  \begin{figure}[H]
    \centering
    \includegraphics[width=\linewidth,height=0.72\textheight,keepaspectratio]{#1}
    \caption{#2}
  \end{figure}}

%=======================================================================
\begin{document}

%------------------------------ Portada --------------------------------
\begin{titlepage}
\centering
\vspace*{3cm}

{\Large\bfseries\color{acento} INFORME DE TALLER PRÁCTICO\par}
\vspace{1cm}

{\Huge\bfseries Terraform + Docker\par}
\vspace{0.8cm}

{\large\itshape Automatizando una arquitectura Node.js + React + Nginx +
MySQL, localmente con Docker\par}

\vspace{2cm}
\rule{\linewidth}{0.4pt}
\vspace{0.6cm}

{\large Jairo Miguel Lara Serrano\par}
\vspace{0.4cm}
{Maestría en Ingeniería de Software\par}
{Desarrollo de Aplicaciones Empresariales\par}
\vspace{0.4cm}
{\today\par}

\vfill
\end{titlepage}

\tableofcontents
\newpage

%=======================================================================
\section{Objetivo del taller}

El taller consiste en desplegar, con Terraform y sin ningún proveedor de
nube, una arquitectura completa de cuatro capas sobre contenedores Docker
locales, y usarla como banco de pruebas para practicar el archivo de
estado (\texttt{terraform.tfstate}).

Los objetivos concretos son cinco:

\begin{enumerate}
  \item Desplegar la arquitectura Nginx (balanceador) $\rightarrow$ React
        (frontend) $\rightarrow$ Node.js (backend, con réplicas)
        $\rightarrow$ MySQL, usando el proveedor \texttt{docker} de
        Terraform.
  \item Ejecutar el flujo completo \texttt{init} $\rightarrow$
        \texttt{plan} $\rightarrow$ \texttt{apply} $\rightarrow$
        \texttt{destroy}.
  \item Practicar de forma guiada el archivo de estado: inspeccionarlo,
        provocar y detectar \emph{drift}, reemplazar recursos puntuales,
        renombrar direcciones sin destruir infraestructura e importar un
        recurso creado a mano.
  \item Escalar el backend cambiando una sola variable y observar cómo
        Terraform recalcula el plan y regenera la configuración de Nginx.
  \item Revisar qué haría falta antes de llevar este mismo patrón a
        producción.
\end{enumerate}

Todo corre en local: no se usó AWS, GCP ni Azure en ningún momento, y el
taller no consumió crédito de nube.

%=======================================================================
\section{Arquitectura desplegada}

Nginx es el único servicio con un puerto publicado al host
(\texttt{8080}). Reparte \texttt{/api/*} entre las réplicas del backend
mediante balanceo \emph{round-robin} y sirve \texttt{/} haciendo proxy al
contenedor del frontend. Todo lo demás vive dentro de la red Docker
privada \texttt{iac-workshop-net}, sin exposición al exterior.

\begin{table}[H]
\centering
\caption{Componentes de la arquitectura}
\begin{tabularx}{\linewidth}{@{}lYl@{}}
\toprule
\textbf{Componente} & \textbf{Función} & \textbf{Puerto} \\
\midrule
Nginx           & Balanceador y \emph{reverse proxy}; único punto de entrada & 8080 (host) \\
React (frontend)& Interfaz servida como estáticos desde su propio Nginx interno & 80 (interno) \\
Node.js/Express & API con \texttt{/api/health} y \texttt{/api/visits}; N réplicas & 3000 (interno) \\
MySQL 8.0       & Base de datos con volumen nombrado para persistencia & 3306 (interno) \\
\bottomrule
\end{tabularx}
\end{table}

Dos detalles del diseño que resultan importantes más adelante:

\begin{itemize}
  \item El backend reporta su \texttt{hostname} en \texttt{/api/health}.
        Como Terraform fija ese hostname explícitamente
        (\texttt{backend-1}, \texttt{backend-2}\ldots), el balanceo de
        carga se puede comprobar visualmente desde el navegador.
  \item La configuración de Nginx no es un fichero estático: se genera en
        el momento del \texttt{apply} con \texttt{templatefile()} a partir
        de la lista de réplicas. Cambiar el número de réplicas cambia ese
        contenido y obliga a recrear el contenedor de Nginx.
\end{itemize}

%=======================================================================
\section{Entorno y estructura del proyecto}

\begin{table}[H]
\centering
\caption{Herramientas utilizadas}
\begin{tabularx}{\linewidth}{@{}lYl@{}}
\toprule
\textbf{Herramienta} & \textbf{Versión mínima} & \textbf{Verificación} \\
\midrule
Docker Desktop / Engine & 24.x  & \texttt{docker version} \\
Terraform CLI           & 1.5.0 & \texttt{terraform -version} \\
Node.js (opcional)      & 20.x  & \texttt{node -v} \\
\bottomrule
\end{tabularx}
\end{table}

El proveedor \texttt{kreuzwerker/docker} habla directamente con el socket
local de Docker, el mismo que usa el comando \texttt{docker}. Por eso no
hacen falta credenciales de ningún tipo. El único acceso a internet que se
necesita es el de la primera ejecución, para descargar el plugin del
proveedor y las imágenes base.

Solo el directorio \texttt{terraform/} contiene código de infraestructura;
\texttt{backend/}, \texttt{frontend/}, \texttt{nginx/} y \texttt{mysql/}
son las aplicaciones y ficheros de configuración que Terraform construye y
despliega.

\begin{shell}
workshop/
  backend/            API Express (health + visits)
  frontend/           UI React que consume la API
  nginx/              plantilla nginx.conf.tpl que Terraform renderiza
  mysql/              esquema inicial de la base de datos
  terraform/
    versions.tf       provider + version de Terraform
    variables.tf      replicas, puerto, credenciales
    locals.tf         lista de alias backend-1, backend-2, ...
    network.tf        red privada iac-workshop-net
    mysql.tf          imagen + volumen + contenedor
    backend.tf        usa "count" para las replicas
    frontend.tf       una sola replica
    nginx.tf          usa templatefile() + upload
    outputs.tf        URLs y nombres de contenedores
\end{shell}

%=======================================================================
\section{Ciclo básico: de \texttt{init} a la aplicación funcionando}

Todos los comandos se ejecutan dentro de \texttt{terraform/}.

\subsection{Paso 1 --- Inicializar}

\begin{shell}
cd workshop/terraform
terraform init
\end{shell}

Terraform descargó el plugin \texttt{kreuzwerker/docker} v4.5.0, preparó el
directorio \texttt{.terraform/} y creó el fichero de bloqueo
\texttt{.terraform.lock.hcl}, que fija esa versión exacta para futuras
ejecuciones.

\captura{01-init}{\texttt{terraform init}: descarga del proveedor y
creación del fichero de bloqueo de versiones.}

\subsection{Paso 2 --- Formatear y validar}

\begin{shell}
terraform fmt
terraform validate
\end{shell}

\texttt{fmt} normaliza el estilo del código y \texttt{validate} revisa la
sintaxis y las referencias internas. Ninguno de los dos necesita que
Docker esté corriendo, así que sirven como comprobación barata antes de
tocar nada real.

\captura{02-fmt-validate}{La configuración es sintácticamente válida.}

\subsection{Paso 3 --- Planificar}

\begin{shell}
terraform plan
\end{shell}

El plan anunció \textbf{11 recursos a crear} y ningún cambio ni
destrucción: la red, el volumen de MySQL, cuatro imágenes y cinco
contenedores. Todavía no se creó nada; este es el momento de leer con
calma qué va a pasar.

\captura{03-plan}{Resumen del plan: \texttt{11 to add, 0 to change, 0 to
destroy}, junto con los valores que tomarán los \emph{outputs}.}

\subsection{Paso 4 --- Aplicar}

\begin{shell}
terraform apply
\end{shell}

\captura{04-apply-confirm}{Terraform pide confirmación explícita antes de
ejecutar el plan; solo se acepta \texttt{yes}.}

Durante el \texttt{apply} se nota que MySQL tarda alrededor de un minuto
en completarse: no basta con que el proceso arranque, porque el recurso
declara un \texttt{healthcheck} con \texttt{wait = true} y Terraform no lo
da por aplicado hasta que MySQL responde \emph{healthy}. El
\texttt{depends\_on} del backend obliga a esperar a que eso ocurra.

\captura{05-apply-complete}{Creación de los 11 recursos en orden de
dependencias. MySQL tarda 1m5s por el \emph{healthcheck}; los backends se
crean después, en paralelo.}

\captura{06-outputs}{\emph{Outputs} finales: URL de la aplicación, URL de
salud de la API y nombres de los contenedores generados.}

\subsection{Paso 5 --- Verificar}

\begin{shell}
docker ps
curl http://localhost:8080/api/health
curl http://localhost:8080/api/visits
\end{shell}

\captura{07-docker-ps}{Los cinco contenedores en ejecución. Solo
\texttt{workshop-nginx} publica un puerto al host
(\texttt{0.0.0.0:8080->80/tcp}); MySQL aparece como \emph{healthy}.}

Abriendo \texttt{http://localhost:8080} se ve la interfaz de React
leyendo el contador de visitas desde MySQL. El botón que dispara seis
peticiones a \texttt{/api/health} muestra el balanceo en vivo: las
respuestas se reparten entre \texttt{backend-1} y \texttt{backend-2}.

\captura{08-app-2-replicas}{Aplicación funcionando con dos réplicas. Las
seis peticiones alternan entre \texttt{backend-2} y \texttt{backend-1}.}

%=======================================================================
\section{El archivo de estado}

Esta es la parte central del taller y la que no tiene equivalente en
Docker Compose. Compose puede levantar los mismos servicios, pero no
mantiene un registro persistente y consultable de qué es cada recurso, con
qué atributos y cómo se relaciona con los demás.

\subsection{Inspeccionar el estado}

\begin{shell}
terraform state list
\end{shell}

Lista las direcciones de todos los recursos administrados, incluyendo el
índice de cada réplica del backend.

\captura{09-state-list}{Las once direcciones del estado. Nótese
\texttt{docker\_container.backend[0]} y
\texttt{docker\_container.backend[1]}: \texttt{count} produce dos recursos
independientes en el grafo.}

\begin{shell}
terraform state show "docker_container.backend[0]"
\end{shell}

\captura{10-state-show-backend}{Todos los atributos conocidos de una
réplica concreta: su ID real de Docker, su IP interna
(\texttt{172.21.0.4}), su hostname (\texttt{backend-1}), su alias de red y
su usuario no root (\texttt{node}). Las variables de entorno aparecen
ocultas por estar marcadas como sensibles.}

\begin{shell}
terraform show
\end{shell}

\captura{11-terraform-show}{Volcado completo del estado en formato
legible: imágenes, red con su subred \texttt{172.21.0.0/16}, volumen con
su punto de montaje real en el host, y los \emph{outputs}.}

\subsection{Idempotencia}

\begin{shell}
terraform plan
\end{shell}

Sin haber tocado nada, Terraform re-consulta la infraestructura real,
compara contra el estado y concluye que no hay nada que hacer. Esto es la
idempotencia declarativa en acción.

\captura{12-plan-sin-cambios}{\texttt{No changes. Your infrastructure
matches the configuration.}}

\subsection{Provocar y detectar \emph{drift}}

El \emph{drift} es la desviación entre la infraestructura real y lo que el
código declara, típicamente por un cambio manual. Se simuló borrando un
contenedor directamente con el CLI de Docker, sin pasar por Terraform:

\begin{shell}
docker rm -f workshop-backend-2
terraform plan
\end{shell}

Terraform detectó que \texttt{docker\_container.backend[1]} ya no existe
en la realidad aunque el estado dice que debería, y propuso recrearlo.

\captura{13-drift-plan}{Detección del \emph{drift}: el recurso desaparecido
se marca para creación, con \texttt{hostname = "backend-2"} y su alias de
red.}

\captura{14-drift-plan-resumen}{Resumen: \texttt{1 to add, 0 to change, 0
to destroy}. Terraform propone tocar solo el recurso afectado.}

\subsection{Corregir el \emph{drift}}

\begin{shell}
terraform apply
\end{shell}

\captura{15-drift-apply}{El \texttt{apply} recrea únicamente el contenedor
faltante.}

\captura{16-drift-apply-resumen}{\texttt{1 added, 0 changed, 0 destroyed}:
el resto del estado seguía coincidiendo con la realidad y no se tocó.}

\subsection{Reemplazar un recurso puntual}

Cuando no hay \emph{drift} pero se quiere forzar una instancia nueva y
limpia de un recurso concreto:

\begin{shell}
terraform apply -replace="docker_container.backend[0]"
\end{shell}

\captura{17-replace-plan}{El plan marca el recurso como \texttt{replaced,
as requested} y usa el símbolo \texttt{-/+} (destruir y luego crear).}

\captura{18-replace-apply}{\texttt{1 added, 0 changed, 1 destroyed}: solo
\texttt{backend-1} se recreó; \texttt{backend-2} quedó intacto.}

Esta bandera sustituye al antiguo \texttt{terraform taint}, desaconsejado
desde Terraform 0.15.2.

\subsection{Renombrar una dirección sin destruir el recurso}

Si se renombra un recurso solo en el código, Terraform planea destruir el
antiguo y crear uno nuevo, aunque en la práctica sea el mismo contenedor.
\texttt{state mv} evita esa recreación innecesaria:

\begin{shell}
terraform state mv docker_container.frontend docker_container.web
\end{shell}

\captura{19-state-mv}{La dirección se mueve dentro del estado:
\texttt{Successfully moved 1 object(s)}.}

\captura{20-state-mv-list}{El estado ya contiene
\texttt{docker\_container.web} en lugar de
\texttt{docker\_container.frontend}, sin que el contenedor real se haya
tocado.}

En este punto el estado y el código quedan momentáneamente
desincronizados, y el siguiente \texttt{plan} falla. Es un error esperado
y didáctico: \texttt{nginx.tf} seguía apuntando a la dirección vieja.

\captura{21-error-referencia}{\texttt{Reference to undeclared resource}:
\texttt{nginx.tf} línea 35 todavía referencia
\texttt{docker\_container.frontend}.}

Tras renombrar también el bloque en el código, el plan vuelve a mostrar
``sin cambios'' en lugar de un \emph{destroy/create}, que es exactamente
el objetivo del ejercicio.

\captura{22-plan-tras-renombrar}{Con código y estado alineados, Terraform
no propone ninguna acción.}

\subsection{Importar un recurso creado manualmente}

Este es probablemente el ejercicio más ilustrativo. Se creó un contenedor
de Redis \textbf{completamente fuera de Terraform}, como si fuera un
recurso heredado:

\begin{shell}
docker network connect iac-workshop-net --alias redis-manual \
  $(docker run -d --name manual-redis redis:7-alpine)
\end{shell}

\captura{23-redis-docker-run}{Docker descarga \texttt{redis:7-alpine} y
crea el contenedor \texttt{manual-redis} al margen de Terraform.}

Después se escribió en \texttt{terraform/redis.tf} el bloque que
\emph{debería} describir ese contenedor, y se adoptó en el estado sin
recrearlo:

\begin{shell}
terraform import docker_container.redis \
  $(docker inspect -f '{{.Id}}' manual-redis)
\end{shell}

\captura{24-terraform-import}{\texttt{Import successful!} El contenedor
pasa a estar administrado por Terraform sin haber sido destruido ni vuelto
a crear.}

\captura{25-import-plan}{El plan posterior muestra
\texttt{docker\_container.redis must be replaced}: algunos atributos del
código no coinciden exactamente con los del contenedor importado (por
ejemplo el \texttt{command}), lo que fuerza una recreación. Es un
resultado normal del ejercicio.}

\captura{26-import-plan-resumen}{\texttt{2 to add, 0 to change, 1 to
destroy}: se crea además \texttt{docker\_image.redis}, porque el recurso
de imagen no soporta \texttt{import} en este proveedor. Como la imagen ya
existía localmente, es un efecto inofensivo.}

\captura{27-import-apply}{El \texttt{apply} completa la adopción del
recurso.}

Docker Compose no tiene ningún comando equivalente a \texttt{terraform
import}: solo puede gestionar lo que él mismo creó. El ejercicio deja
además una lección extra: no todos los tipos de recurso de un proveedor
soportan importación, y conviene revisar la documentación antes de
asumirlo.

\subsection{Limpieza del ejercicio y una observación}

\begin{shell}
rm terraform/redis.tf
docker rm -f manual-redis
terraform state rm docker_container.redis docker_image.redis
\end{shell}

\captura{28-import-limpieza}{\texttt{terraform state rm} saca los dos
recursos del estado sin destruirlos en la realidad. Nótese, sin embargo,
el mensaje \texttt{rm: terraform/redis.tf: No such file or directory}.}

Merece la pena detenerse en ese detalle. El comando se ejecutó desde
dentro de \texttt{terraform/}, así que la ruta relativa
\texttt{terraform/redis.tf} no resolvía y el fichero \textbf{no se borró}.
El resultado fue que, en el siguiente \texttt{apply} (durante el escalado
de la sección siguiente), Terraform volvió a crear Redis a partir de ese
\texttt{redis.tf} superviviente.

Lejos de ser solo un descuido, es una demostración práctica del principio
de fondo del taller: \textbf{el código es la fuente de verdad}. Sacar algo
del estado no lo elimina de la configuración, y mientras el bloque siga
declarado Terraform lo reconstruirá. La ruta correcta desde
\texttt{terraform/} habría sido simplemente \texttt{rm redis.tf}.

\subsection{Resumen de la comparación con Docker Compose}

\begin{table}[H]
\centering
\caption{Qué puede hacer cada herramienta}
\begin{tabularx}{\linewidth}{@{}Yll@{}}
\toprule
\textbf{Capacidad} & \textbf{Compose} & \textbf{Terraform} \\
\midrule
Detectar si la configuración cambió              & Sí       & Sí \\
Detectar \emph{drift} por cambios manuales       & No       & Sí (\texttt{plan}) \\
Vista previa con \emph{diff} explícito           & No       & Sí (\texttt{plan}) \\
Consultar atributos completos de un recurso      & Limitado & Sí (\texttt{state show}) \\
Renombrar sin recrear el recurso real            & No       & Sí (\texttt{state mv}) \\
Adoptar un recurso creado por fuera              & No       & Sí (\texttt{import}) \\
Reemplazar un recurso puntual bajo demanda       & No       & Sí (\texttt{-replace}) \\
Bloqueo contra aplicaciones concurrentes         & No       & Sí (\emph{backend} remoto) \\
\bottomrule
\end{tabularx}
\end{table}

%=======================================================================
\section{Escalar el backend}

Basta cambiar una sola variable en \texttt{terraform.tfvars}:

\begin{shell}
backend_replica_count = 4
\end{shell}

\begin{shell}
terraform plan
\end{shell}

\captura{29-escalar-plan}{El plan añade \texttt{workshop-backend-3} y
\texttt{workshop-backend-4} a la lista de \emph{outputs}.}

Lo interesante no son las réplicas nuevas, sino que Nginx también aparece
marcado para recreación. Su configuración se genera con
\texttt{templatefile()} a partir de la lista de réplicas; al cambiar esa
lista cambia el contenido del fichero, y como el fichero se sube al crear
el contenedor, no se puede parchear en caliente. Es el principio de
infraestructura inmutable, visible en la práctica.

\captura{30-escalar-nginx-replace}{\texttt{docker\_container.nginx must be
replaced}: la recreación no viene de un cambio manual, sino de una
dependencia sobre datos que cambiaron.}

\begin{shell}
terraform apply
\end{shell}

\captura{31-escalar-apply}{\texttt{5 added, 0 changed, 1 destroyed}: dos
backends nuevos, Nginx destruido y recreado, y Redis recreado a partir del
\texttt{redis.tf} que había sobrevivido a la limpieza.}

Volviendo al navegador, el botón de balanceo ahora reparte las peticiones
entre cuatro hostnames distintos.

\captura{32-app-4-replicas}{La misma interfaz, ahora repartiendo entre
\texttt{backend-1} a \texttt{backend-4}. El contador de visitas subió a 2,
leído desde el volumen persistente de MySQL.}

%=======================================================================
\section{Destruir el entorno}

\begin{shell}
terraform destroy
\end{shell}

\captura{33-destroy}{\texttt{Destroy complete! Resources: 15 destroyed.}
Terraform elimina todo en orden inverso de dependencias: primero los
contenedores, luego las imágenes, y al final el volumen y la red.}

Los 15 recursos son los 11 originales más los 2 backends añadidos al
escalar y los 2 de Redis. Conviene subrayar que \texttt{destroy} se llevó
también el volumen \texttt{iac-workshop-mysql-data} y, con él, los datos
de MySQL. En un entorno con datos reales, \texttt{lifecycle \{
prevent\_destroy = true \}} sobre ese volumen es la red de seguridad
contra exactamente este comando.

%=======================================================================
\section{Consideraciones antes de llevarlo a producción}

El proveedor \texttt{docker} es excelente para aprender y para desarrollo
local, pero no está pensado para producción. Los ajustes que un equipo
debería hacer antes de dar ese salto:

\begin{enumerate}
  \item \textbf{Cambiar de orquestador.} Este proveedor administra un solo
        host: no reinicia nada si el host cae, no distribuye carga entre
        máquinas y no tiene autoescalado. El mismo enfoque declarativo se
        aplica con Terraform sobre Kubernetes o ECS.
  \item \textbf{Backend remoto de estado, con bloqueo.} El
        \texttt{terraform.tfstate} local es inviable en equipo: dos
        personas aplicando a la vez pueden corromperlo.
  \item \textbf{Nunca variables sensibles con \texttt{default}.} Las
        contraseñas no deberían tener valor por defecto, ni vivir en un
        \texttt{.tfvars} versionado, sino inyectarse desde un gestor de
        secretos.
  \item \textbf{Tags de imagen inmutables, nunca \texttt{:latest}}, para
        poder hacer \emph{rollback} con certeza.
  \item \textbf{Límites de recursos por contenedor}
        (\texttt{memory}, \texttt{cpu\_shares}).
  \item \textbf{TLS en el punto de entrada}, terminando en el balanceador.
  \item \textbf{Backups del dato}, o mejor, una base de datos gestionada.
  \item \textbf{\texttt{prevent\_destroy}} en los recursos con estado
        crítico.
  \item \textbf{Pipeline de CI/CD con aprobación humana}, políticas como
        código y escaneo de imágenes.
  \item \textbf{Usuario no root en los contenedores} (ya se cumple en el
        backend, que usa \texttt{USER node}).
  \item \textbf{Monitoreo y logging centralizados.}
  \item \textbf{Ambientes separados con estados completamente distintos},
        para limitar el radio de explosión de un error.
\end{enumerate}

%=======================================================================
\section{Referencia rápida de comandos}

\begin{table}[H]
\centering
\caption{Comandos utilizados en el taller}
\begin{tabularx}{\linewidth}{@{}lY@{}}
\toprule
\textbf{Comando} & \textbf{Qué hace} \\
\midrule
\texttt{terraform init}    & Descarga proveedores y prepara el directorio \\
\texttt{terraform fmt}     & Normaliza el formato del código \\
\texttt{terraform validate}& Revisa sintaxis y referencias, sin tocar Docker \\
\texttt{terraform plan}    & Calcula y muestra el plan, sin aplicar nada \\
\texttt{terraform apply}   & Ejecuta el plan (pide confirmación) \\
\texttt{terraform destroy} & Elimina todos los recursos administrados \\
\texttt{terraform state list} & Lista las direcciones del estado \\
\texttt{terraform state show} & Muestra los atributos de un recurso \\
\texttt{terraform show}    & Vuelca el estado completo en formato legible \\
\texttt{terraform state mv}& Renombra una dirección sin destruir el recurso \\
\texttt{terraform state rm}& Deja de administrar un recurso, sin destruirlo \\
\texttt{terraform import}  & Adopta un recurso ya existente \\
\texttt{terraform apply -replace} & Fuerza la recreación de un recurso \\
\texttt{terraform output}  & Muestra los valores de salida \\
\bottomrule
\end{tabularx}
\end{table}

%=======================================================================
\section{Conclusiones}

\begin{itemize}
  \item El flujo \texttt{init} $\rightarrow$ \texttt{plan} $\rightarrow$
        \texttt{apply} se completó sin incidencias: 11 recursos creados y
        la aplicación respondiendo en \texttt{localhost:8080}, con el
        balanceo entre réplicas verificable desde el navegador.
  \item El archivo de estado es la diferencia real frente a Docker
        Compose. Detectar el borrado manual de un contenedor, renombrar
        una dirección sin recrear nada y adoptar un contenedor creado a
        mano son operaciones que Compose sencillamente no puede ofrecer,
        porque no mantiene un registro consultable de la infraestructura.
  \item El escalado dejó ver el efecto en cascada más interesante del
        taller: cambiar una variable no solo añadió réplicas, sino que
        obligó a recrear Nginx, porque su configuración se deriva de esa
        misma variable. Terraform razona sobre el grafo de dependencias,
        no sobre ficheros sueltos.
  \item El descuido en la limpieza de Redis resultó ser el ejemplo más
        claro de todo el ejercicio: sacar un recurso del estado no lo
        borra del código, y mientras el bloque siga declarado Terraform lo
        reconstruirá en el siguiente \texttt{apply}.
  \item El \texttt{destroy} final eliminó los 15 recursos, incluido el
        volumen de MySQL, confirmando por qué \texttt{prevent\_destroy} es
        necesario en cuanto hay datos que importan.
\end{itemize}

\end{document}
